\begin{pastexam}[label={px:ntut-110-prog-02}]{北科大資工所 110 學年度 第 2 題}
（程式追蹤與填空）

\textbf{Problem 2 [12\%] [ 2-1, 3\%] [ 2-2-i, 2-2-ii, each 3\%] [ 2-3-i, 1\%, 2-3-ii, 2\%]}

The following Python code will print the Fibonacci sequence. The output of f(8) is ``0, 1, 1, 2, 3, 5, 8, 13,''.
\begin{minted}[linenos, numbersep=1.5mm, xleftmargin=2em, fontsize=\small]{python}
def fibonacci(n):
    if n<=0:
        return 0
    elif n==1 or n==2:
        return 1
    else:
        return fibonacci(n-1)+fibonacci(n-2)

def f(x):
    i = 0
    while i<x:
        print(fibonacci(i), end=', ')
        i = i + 1
\end{minted}
\begin{enumerate}[label=(\arabic*)]
\item (3\%) If f(8) is executed, how many times will fibonacci(4) be called?
\item (6\%) Please complete the following ``Loop version'' Python code of the Fibonacci function. The output of h(8) is ``0, 1, 1, 2, 3, 5, 8, 13,''.\\
(i) (3\%) Line 4: \underline{\hspace{4em}}\\
(ii) (3\%) Line 5: \underline{\hspace{4em}}
\begin{minted}[linenos, numbersep=1.5mm, xleftmargin=2em, fontsize=\small]{python}
def fibonaccii(n):
    data = [0 for i in range(0, n+1)]
    data[1]=1
    for i in range(_____, n+1):        # (i)
        data[i] = _______________      # (ii)
    return data

def h(x):
    data = fibonaccii(x)
    for i in range(x):
        print(data[i], end=', ')
\end{minted}
\item (3\%) The following code is modified from the code of the above problem (1).
\begin{minted}[linenos, numbersep=1.5mm, xleftmargin=2em, fontsize=\small]{python}
def fibonacci3(n, data):
    if n>5: print(n,' * ', end='')
    if n<=0: return 0
    elif n==1 or n==2: data[n] = 1
    elif data[n]==0:
        data[n] = fibonacci3(n-1, data)+fibonacci3(n-2, data)
    return data[n]

def p(x):
    data = [0 for i in range(x+1)]
    for i in range(x):
        fibonacci3(i, data)
\end{minted}
(i) (1\%) What is the output of p(8)?\\
(ii) (2\%) Please compare fibonacci3 with Fibonacci function according to the computing time and space.
\end{enumerate}
\end{pastexam}

\begin{pastsoln}
答案：(1) $7$ 次；(2)(i) \texttt{2}，(ii) \texttt{data[i-1]+data[i-2]}；(3)(i) 依序印出 $6$、$7$、$6$，空格的寫法見下；(ii) 見下。

(1) \texttt{f(8)} 依序呼叫 \texttt{fibonacci(0)} 到 \texttt{fibonacci(7)}，各自是一棵獨立的遞迴樹。
令 $C(n)$ 為呼叫一次 \texttt{fibonacci(n)} 時，\texttt{fibonacci(4)} 被呼叫的次數（含它自己）。
$n \le 3$ 時遞迴樹裡只有 $3$ 以下的參數，$C(n) = 0$；$C(4) = 1$；$n \ge 5$ 時第 7 行呼叫
\texttt{n-1} 和 \texttt{n-2}，所以 $C(n) = C(n - 1) + C(n - 2)$：
\[
C(5) = C(4) + C(3) = 1, \quad C(6) = C(5) + C(4) = 2, \quad C(7) = C(6) + C(5) = 3 .
\]
合計 $C(0) + C(1) + \cdots + C(7) = 0 + 0 + 0 + 0 + 1 + 1 + 2 + 3 = 7$。 這正是資結書第~18.1.1~
節說的重疊子問題：同一個參數被重算很多次（\texttt{f(8)} 一共呼叫 \texttt{fibonacci} $60$ 次，
其中 \texttt{fibonacci(4)} 佔 $7$ 次）。

(2) \texttt{data[0]} $= 0$、\texttt{data[1]} $= 1$ 已經設好，迴圈要從 $2$ 開始：(i) 填 \texttt{2}，
\texttt{range(2, n+1)} 讓 \texttt{i} 跑 $2, 3, \ldots, n$。 (ii) 填 \texttt{data[i-1]+data[i-2]}。
這是資結書第~18.1.3~節的由下而上：算 \texttt{data[i]} 時，用到的兩格都已經填好。
\texttt{h(8)} 呼叫 \texttt{fibonaccii(8)} 得到 \texttt{data[0..8]}，印出前 $8$ 個，是
\texttt{0, 1, 1, 2, 3, 5, 8, 13,}，和題目要的一樣。

(3)(i) \texttt{p(8)} 對 $i = 0, 1, \ldots, 7$ 依序呼叫 \texttt{fibonacci3(i, data)}，
\texttt{data} 是同一個串列，前面算過的值會留著。 第 2 行只在參數大於 $5$ 時印東西，而且它在檢
查 \texttt{data[n]} 之前就印，所以每一次參數是 $6$ 或 $7$ 的呼叫都會印一次：
\begin{itemize}
\item $i \le 5$：參數都不超過 $5$，不印。 $i = 5$ 結束時 \texttt{data[1..5]} 都已經算好。
\item $i = 6$：印出 \texttt{6}；\texttt{data[6]} 還是 $0$，第 6 行呼叫 \texttt{fibonacci3(5)}、
\texttt{fibonacci3(4)}，兩個都直接回傳，不印。
\item $i = 7$：印出 \texttt{7}；\texttt{data[7]} 是 $0$，第 6 行先呼叫 \texttt{fibonacci3(6)}，
它在第 2 行又印出 \texttt{6}，再回傳 \texttt{data[6]}；接著 \texttt{fibonacci3(5)} 不印。
\end{itemize}
所以依序印出 6、7、6 三次。 Python 的 \texttt{print(n, ' * ', end='')} 會在兩個參數之間補一個空
格，每次印出的是「數字、兩個空格、星號、一個空格」，完整輸出是
\begin{center}
\verb*|6  * 7  * 6  * |
\end{center}
（\verb*| | 表示空格。 只寫 \texttt{6 * 7 * 6 *} 的話，三次的順序對了，空格數不同。）

(ii) \texttt{fibonacci} 沒有記錄算過的值，呼叫次數 $T(n) = 1 + T(n - 1) + T(n - 2)$（$n \ge 3$），
和費氏數一樣以約 $1.618^n$ 的速度成長，時間是指數級的；額外空間只有遞迴的堆疊，深度
$\Theta(n)$。 \texttt{fibonacci3} 是資結書第~18.1.2~節的記憶化：每個 \texttt{data[k]} 第一次
被問到時算一次，之後直接回傳，所以算 \texttt{fibonacci3(n)} 只要 $O(n)$ 次呼叫，時間 $O(n)$；
\texttt{p(8)} 因為前面的值都留著，每個 $i$ 最多再多兩次呼叫，總共 $18$ 次。 代價是多一個長度
$n + 1$ 的串列 \texttt{data}，空間 $\Theta(n)$（遞迴深度一樣最多 $\Theta(n)$）。 也就是用
$\Theta(n)$ 的記憶體，把指數時間換成線性時間。
\end{pastsoln}
